% open-logic-complete-pt.tex
% master file to produce a complete release version
% without references, markup, etc.

\documentclass[openany]{memoir}

% we're compiling from main directory
\newcommand{\olpath}{../..}
\newcommand*{\ollangid}{pt}
\newcommand*{\ollanguage}{brazilian}

\usepackage{mathpazo}
\usepackage{helvet}
\usepackage{courier}
\linespread{1.05} % Palatino looks better with this

\usepackage{gitinfo2}

\input{\olpath/sty/open-logic.sty}
\input{\olpath/sty/open-logic-defer.sty}

\includeenv{editorial}

\problemsperchapter

\let\cleardoublepage\clearpage

\pagestyle{plain}

\begin{document}

\begin{titlingpage}
\begin{raggedleft}
\HUGE
\selectfont\bfseries\sffamily
THE OPEN LOGIC TEXT
\vskip 3ex
\normalfont\huge
Complete Build\\
PORTUGUESE
\vskip 4ex
\textbf{\href{http://openlogicproject.org/}{Open Logic Project}}
\vskip 4ex
\Large
Revision\gitRel: \gitAbbrevHash{} (\gitBranch)\\
\gitAuthorDate

\end{raggedleft}

\vfill

\ollicense

\end{titlingpage}

\pagestyle{giheadings}

\chapter*{Sobre o Open Logic Project}
\addcontentsline{toc}{chapter}{Sobre o Open Logic Project}


O \textit{Open Logic Text} é um livro-texto colaborativo de código aberto
de meta-lógica formal e métodos formais, começando em um nível intermediário
(ou seja, após um curso introdutório de lógica formal). Embora voltado a um
público não matemático (em particular, estudantes de filosofia e
ciência da computação), é rigoroso.

A cobertura de alguns tópicos atualmente incluídos pode
ainda não estar completa, e muitas seções ainda exigem revisão substancial.
Planejamos expandir o texto para cobrir mais tópicos no
futuro. Também planejamos adicionar recursos ao texto, como um
glossário, uma lista de leituras complementares, notas históricas, imagens, melhores
explicações, seções que explicam a relevância dos resultados para
filosofia, ciência da computação e matemática, e mais problemas e
exemplos. Se você encontrar um erro ou tiver uma sugestão,
\href{https://github.com/OpenLogicProject/OpenLogic/wiki/Contributing}{informe a equipe do projeto}.

O projeto opera no espírito do código aberto. O texto não só está
disponível gratuitamente; fornecemos o código-fonte LaTeX sob a
licença Creative Commons Atribuição, que dá a qualquer pessoa o direito de
baixar, usar, modificar, reorganizar, converter e redistribuir nosso
trabalho, desde que deem o devido crédito.
Consulte o site do Open Logic Project em
\href{http://openlogicproject.org/}{openlogicproject.org} para
informações adicionais.


\tableofcontents*

\olimport[content]{content}

\photocredits

\bibliographystyle{\olpath/bib/natbib-oup}
\bibliography{\olpath/bib/open-logic}

\end{document}

